%=======================================================================
%  LABORATORIO N.o 7 - AMPLIFICADORES CON TRANSISTOR BIPOLAR
%  Curso : Circuitos Electronicos  (2026-2)
%  UCSP  - Facultad de Ingenieria - Ing. Electronica y de Telecomunicaciones
%  Docente     : Ebert San Roman Castillo
%  Laboratorio : Ing. Jeanpier Ancori Sanchez
%
%  Version 2 - Practica de DOS SEMANAS con las tres configuraciones del BJT
%  y la etapa de potencia (la v1 tenia tres experimentos, sin base comun).
%  Semana 1: experimentos 1 y 2 (amplificadores de voltaje).
%  Semana 2: experimentos 3 y 4 (corriente y potencia).
%    1. Amplificador de VOLTAJE: emisor comun. Es el circuito 3 (punto
%       medio) del Laboratorio 5 con Cout y RL; el alumno ya lo tiene armado
%       y polarizado. Se mide Av, Zin, Ai y Ap con RL = 10 k y con RL = 100.
%    2. Amplificador de VOLTAJE: base comun, sobre la misma placa del exp. 1
%       (se quita CE, la base va a tierra en AC con CB = 100 uF y la senal
%       entra por el emisor). Mismo punto Q: Av ~ 10 veces mayor que el
%       emisor comun pero Zin ~ 11 Ohm y Ai < 1.
%    3. Amplificador de CORRIENTE: colector comun (seguidor de emisor) con
%       RL = 100. Primero la carga sin transistor (vo ~ 4 mV), luego con el
%       seguidor y despues detras del emisor comun: la ganancia en 100 Ohm
%       pasa de ~0,9 a ~15.
%    4. Amplificador de POTENCIA: push-pull complementario 2N2222A/2N2907A,
%       primero en clase B (distorsion de cruce) y luego en clase AB con dos
%       1N4148. Pout, PDC, rendimiento, Ap.
%  Mejoras de la v2 en el marco teorico: "tierra para la senal" (el emisor
%  visto en DC y en AC, y por que VCC es tierra en AC), el nombre de cada
%  configuracion, la base comun con su circuito y la tabla comparativa con
%  cuatro columnas.
%  Formato del Laboratorio 6 (Datos -> Calcule -> Arme y mida -> Anote,
%  cuadro resumen, cuestionario de 4 preguntas (una por experimento), conclusiones 5 + 5); no se
%  imprimen los resultados. Marco teorico ampliado a pedido: el alumno no ha
%  visto aun el modelo de pequena senal, asi que se dan las formulas con
%  r_e = 26 mV/I_E y un ejemplo resuelto.
%  Zin se mide con una Rs en serie con W1: Zin = Rs vin/(vs - vin).
%  Instrumental: Analog Discovery Studio. Exp. 1, 2 y 3 con el riel de +12 V
%  (0,2 A; consumo < 10 mA). Exp. 4 con las fuentes programables V+ = 5 V y
%  V- = -5 V (700 mA); corriente pico en RL ~ 50 mA. W1 entrega como maximo
%  ~5 mA pico en el exp. 4 (Zin ~ 480 Ohm) y ~0,45 mA en el exp. 2
%  (Zin ~ 11 Ohm, con Rs = 100); los demas, menos de 0,1 mA.
%
%  VALORES DE REFERENCIA PARA EL DOCENTE (beta = 200, VBE = 0,7 V,
%  r_e = 26 mV/IE; no se imprimen):
%    Exp. 1: VTH = 1,83 V, RTH = 1,53 k, IE = 2,27 mA, re = 11,5 Ohm,
%      VB/VE/VC = 1,83/1,13/7,0 V.
%      Av sin RL = -2200/111,5 = -19,7.
%      Zin = 10k || 1,8k || 200*111,5 = 1,43 k.
%      RL = 10 k : Av = -1803/111,5 = -16,2 ; Ai = 16,2*1,43/10 = 2,3 ;
%                  Ap = 37 (15,7 dB).
%      RL = 100  : Av = -95,7/111,5 = -0,86 ; Ai = 12,3 ; Ap = 10,6 (10 dB).
%      Con W1 = 50 mV pico y Rs = 1 k: vin ~ 30 mV pico, vo ~ 0,48 V pico.
%    Exp. 2 (base comun, mismo punto Q, RE = 490 sin desacoplar, Rs = 100,
%      W1 = 50 mV pico): Zin = 490 || 11,5 = 11,2 Ohm ;
%      vin = 50 mV * 11,2/111,2 = 5,0 mV pico.
%      Av sin RL = 2200/11,5 = 191 ; vo = 0,96 V pico ;
%      Av(vs -> vo) = 2200/(100 + 11,5) = 19,7 (igual que el emisor comun).
%      RL = 10 k : Av = 1803/11,5 = 157 ; Ai = 157*11,2/10k = 0,18
%                  (= alfa RC/(RC + RL)) ; Ap = 28 (14,5 dB).
%      Con 0,2 V pico en W1 la salida se deforma (ie ~ 0,8 IE).
%    Exp. 3: VTH = 6 V, RTH = 23,5 k, IE = 5,3/(1k + 23,5k/200) = 4,75 mA,
%      re = 5,5 Ohm, VE = 4,75 V, P_RE = 23 mW.
%      RE || RL = 90,9 ; Av = 90,9/96,4 = 0,943 ;
%      Zin = 23,5k || 200*96,4 = 23,5k || 19,3k = 10,6 k ;
%      Ai = 0,943*10,6k/100 = 100 ; Ap = 94 (19,8 dB).
%      Recorte del semiciclo negativo en vo = -IE (RE||RL) = -0,43 V.
%      Cascada (CE sin RL -> seguidor -> 100 Ohm), vin = 20 mV pico:
%      Av1 = -(2,2k || 10,6k)/111,5 = -16,3 ; Av = -16,3*0,943 = -15,4
%      (contra -0,86 con el CE directo a 100 Ohm).
%      Sin transistor (RL = 100 detras de Rs = 10 k): vo = 0,4*100/10,1k
%      = 4 mV pico, contra ~0,19 V con el seguidor.
%    Exp. 4 (+-5 V, RL = 47 Ohm, W1 = 3 V pico, Rs = 100):
%      Clase B: Vp ~ 3 - 0,7 = 2,3 V ; Pout = 2,3^2/94 = 56 mW ;
%        PDC = 2*5*2,3/(pi*47) = 156 mW ; eta = (pi/4)(2,3/5) = 36 %.
%        PD(max) por transistor = 25/(pi^2*47) = 54 mW.
%      Clase AB: I por los diodos = (10 - 1,3)/2 k = 4,35 mA (43 mW extra);
%        Zin ~ RB1 || RB2 = 500 Ohm ; Vp ~ 2,4 V ; Pout ~ 61 mW ;
%        PDC ~ 156 + 43 = 200 mW ; eta ~ 30 % ; Av ~ 0,95.
%    Cuestionario 1: sin CE, Av = -2200/(11,5 + 490) = -4,4 ; Zin sube a
%      10k || 1,8k || 100k = 1,51 k.
%    Cuestionario 2 (base comun): con Rs = 1 k, Av(vs -> vo) =
%      2200/(1000 + 11,5) = 2,2 ; con Rs = 0 seria 191.
%    Cuestionario 3: vo max sin recorte = 0,43 V pico ; con RE = 470 Ohm,
%      IE ~ 10 mA, vo max ~ 10 mA * 82,5 = 0,83 V pico (P_RE = 47 mW).
%    Cuestionario 4: con Vp = 4,3 V, eta = (pi/4)(4,3/5) = 67,5 % ;
%      Pout = 4,3^2/94 = 197 mW.
%
%  SIMULACION (ngspice 26 de Fusion/Eagle; modelos BC548B de NXP, 2N2222A,
%  2N2907A y 1N4148; netlists y grafica en ../Simulaciones):
%    Exp. 1: VB/VE/VC = 1,82/1,16/6,82 V, IC = 2,36 mA (beta ~ 303).
%      Av sin RL = -19,5 ; Zin = 1,46 k.
%      RL = 10 k : Av = -16,0 ; Ai = 2,33 ; Ap = 37,2.
%      RL = 100  : Av = -0,85 ; Ai = 12,5 ; Ap = 10,6.
%    Exp. 2: VB/VE/VC = 1,82/1,16/6,82 V (el mismo punto Q del exp. 1).
%      Zin = 12,4 Ohm ; vin = 5,5 mV pico ; Av sin RL = 173 ;
%      Av(vs -> vo) = 19,0.
%      RL = 10 k : Av = 144 ; Ai = 0,175 ; Ap = 25.
%      Con 0,2 V pico en W1: vo +2,96/-3,15 V (asimetrica, se deforma).
%    Exp. 3: VB = 5,39 V, VE = 4,71 V, IC = 4,68 mA (beta ~ 181).
%      Av = 0,937 ; Zin = 10,4 k ; Ai = 97,8 ; Ap = 91,7.
%      Con 1,2 V pico en W1 el minimo de vo se aplana en -0,43 V.
%      Cascada: Av1 = -16,1 ; Av total = -15,0.
%    Exp. 4 clase B : vo pico 2,22 V ; Av(rms) = 0,68 (cruce) ;
%      Pout = 43,7 mW ; PDC = 126 mW ; eta = 34,6 %.
%    Exp. 4 clase AB: VB1/VB2 = +0,65/-0,65 V, Iq = 0,6 mA ;
%      vo +2,32/-2,28 V ; Av(rms) = 0,92 ; Zin = 483 Ohm ;
%      Pout = 55,4 mW ; Pin = 6,4 mW ; Ap = 8,7 ; PDC = 196 mW ;
%      eta = 28 %.
%=======================================================================

\documentclass[11pt,a4paper]{article}

%------------------------------------------------------------- idioma
% Compilar con XeLaTeX (no con pdflatex): Trebuchet MS es una fuente TrueType
% del sistema y solo fontspec, bajo XeLaTeX o LuaLaTeX, sabe cargarla.
\usepackage{fontspec}
\setmainfont{Trebuchet MS}[
  Path            = C:/Windows/Fonts/,
  Extension       = .ttf,
  UprightFont     = trebuc,
  BoldFont        = trebucbd,
  ItalicFont      = trebucit,
  BoldItalicFont  = trebucbi,
  Ligatures       = TeX
]
\usepackage[spanish,es-nodecimaldot]{babel}
\usepackage{microtype}

%----------------------------------------------------------- geometria
\usepackage[left=2.4cm,right=2.4cm,top=3.1cm,bottom=2.6cm,headheight=30pt,
            footskip=1.2cm]{geometry}

%------------------------------------------------------------ graficos
\usepackage{graphicx}
\usepackage{xcolor}
\usepackage{float}
\usepackage{placeins}
\usepackage[font=small,labelfont=bf,labelsep=period]{caption}
\usepackage{subcaption}
\graphicspath{{imagenes/}}
\usepackage[american,siunitx]{circuitikz}
% babel-spanish activa < > " y rompe las opciones de TikZ sin esta libreria.
\usetikzlibrary{calc,patterns,babel,arrows.meta,positioning}

%-------------------------------------------------------------- tablas
\usepackage{array}
\usepackage{tabularx}
\usepackage{multirow}

%----------------------------------------------------- matematicas / SI
\usepackage{amsmath,amssymb}
\usepackage{siunitx}
\sisetup{
  output-decimal-marker = {,},
  per-mode              = symbol,
  mode                  = match,
  detect-weight         = true,
  detect-family         = true
}

%--------------------------------------------------------------- listas
\usepackage{enumitem}
\setlist[itemize]{itemsep=3pt,topsep=4pt,parsep=0pt}
\setlist[enumerate]{itemsep=8pt,topsep=4pt,parsep=0pt}

%------------------------------------------------------------ comandos
\usepackage{pgffor}

% Lineas punteadas para respuestas del alumno: \lineasresp{n}
\newcommand{\lineasresp}[1]{%
  \par\vspace{6pt}%
  \foreach \i in {1,...,#1}{\noindent\rule{\linewidth}{0.4pt}\par\vspace{8pt}}%
}

% Recuadro "Figura / Comentario": una celda alta para pegar la captura del
% osciloscopio y otra debajo para el comentario.
\newcommand{\cuadrorespuesta}{%
  \par\vspace{6pt}\noindent
  \begin{tabularx}{\linewidth}{|>{\centering\arraybackslash}X|}
    \hline
    \rule{0pt}{4.4cm}\raisebox{2.1cm}{\textcolor{black!45}{Figura}}\\
    \hline
    \rule{0pt}{2.0cm}\raisebox{0.9cm}{\textcolor{black!45}{Comentario}}\\
    \hline
  \end{tabularx}\par\vspace{8pt}%
}

% Celda vacia de altura comoda para escribir a mano.
\newcommand{\hueco}{\rule{0pt}{0.5cm}}

% Paralelo de resistencias: R_1 \pll R_2
\newcommand{\pll}{\mathbin{\|}}

%----------------------------------------------------- compaginacion
\usepackage{needspace}
% \needspace mide mal el espacio libre despues de figuras y tablas [H] y deja
% paginas medio vacias; \Needspace es la variante exacta del mismo paquete.
\let\needspace\Needspace
\usepackage{etoolbox}
\pretocmd{\section}{\needspace{5\baselineskip}}{}{}
\pretocmd{\subsection}{\needspace{4\baselineskip}}{}{}
\widowpenalty=10000
\clubpenalty=10000

%---------------------------------------------------- encabezado / pie
\usepackage{fancyhdr}
\pagestyle{fancy}
\fancyhf{}
\fancyhead[L]{\footnotesize UCSP --- Facultad de Ingeniería\\
              Ing. Electrónica y de Telecomunicaciones}
\fancyhead[R]{\footnotesize 2026-2 Circuitos Electrónicos\\
              Ebert San Román Castillo}
\fancyfoot[L]{\footnotesize Laboratorio N.\textsuperscript{o} 7 --- Amplificadores con transistor bipolar}
\fancyfoot[R]{\footnotesize Página \thepage\ de \pageref{LastPage}}
\renewcommand{\headrulewidth}{0.4pt}
\renewcommand{\footrulewidth}{0.4pt}
\usepackage{lastpage}

%------------------------------------------------------- numeracion
\renewcommand{\thesection}{\Roman{section}}
\renewcommand{\thesubsection}{\arabic{subsection}}
\renewcommand{\thesubsubsection}{\thesubsection.\arabic{subsubsection}}
\setcounter{secnumdepth}{3}

%-------------------------------------------------------- hiperenlaces
\usepackage[hidelinks,pdfusetitle]{hyperref}
\usepackage{xurl}
\hypersetup{
  pdftitle={Laboratorio 7 - Amplificadores con transistor bipolar},
  pdfsubject={Circuitos Electronicos - Guia de practicas},
  pdfauthor={Ebert San Roman Castillo; Jeanpier Ancori Sanchez},
  pdfkeywords={electronica, transistor, BJT, amplificador, ganancia de voltaje,
               ganancia de corriente, ganancia de potencia, emisor comun,
               colector comun, seguidor de emisor, push-pull, clase B, clase AB,
               distorsion de cruce, rendimiento, impedancia de entrada,
               Analog Discovery Studio, Proteus, Multisim, osciloscopio}
}

% babel-spanish rotula las tablas como "Cuadro"; los originales dicen "Tabla".
\addto\captionsspanish{\renewcommand{\tablename}{Tabla}}

\begin{document}

%=======================================================================
%  CABECERA DE LA GUIA
%=======================================================================
\begin{center}
    {\normalsize\bfseries SEGUNDA UNIDAD: TRANSISTOR DE UNIÓN BIPOLAR (BJT)}
\end{center}

\vspace{0.55cm}

\begingroup
\renewcommand{\tabularxcolumn}[1]{m{#1}}
\renewcommand{\arraystretch}{2.0}
\setlength{\tabcolsep}{12pt}
\noindent
\begin{tabularx}{\textwidth}{|>{\centering\arraybackslash}m{2.2cm}|>{\centering\arraybackslash}X|>{\centering\arraybackslash}m{3.0cm}|}
\hline
\rule{0pt}{2.1cm}{\fontsize{58}{58}\selectfont\bfseries 7}
    & \begin{minipage}[c]{\linewidth}\centering
          {\large\bfseries Guía de Prácticas}\\[7pt]
          {\Large\bfseries Amplificadores con\\ Transistor Bipolar}
      \end{minipage}
    & {\Large\bfseries Nota}\\
\hline
\multicolumn{3}{|l|}{\rule{0pt}{1.3cm}\large\textbf{Grupo:}~\rule{4.6cm}{0.4pt}%
    \hfill\large\textbf{Fecha:}~\rule{4.6cm}{0.4pt}\hspace{2mm}}\\
\hline
\multicolumn{3}{|l|}{\rule{0pt}{1.2cm}\large\textbf{Alumno(s):}~\rule{11.2cm}{0.4pt}\hspace{2mm}}\\
\multicolumn{3}{|l|}{\rule{0pt}{1.2cm}\phantom{\large\textbf{Alumno(s):}}~\rule{11.2cm}{0.4pt}\hspace{2mm}}\\
\hline
\end{tabularx}
\endgroup

\vspace{0.8cm}

% Encabezado de cada paso dentro de un experimento: \paso{1}{Datos}
\newcommand{\paso}[2]{%
  \par\needspace{10\baselineskip}\vspace{8pt}\noindent
  \colorbox{black!8}{\makebox[\dimexpr\linewidth-2\fboxsep\relax][l]{%
    \textbf{Paso #1 \,·\, #2}}}\par\vspace{4pt}}

%=======================================================================
\section{Objetivos}
%=======================================================================
\begin{itemize}
    \item Calcular, simular y medir las ganancias de voltaje, de corriente y de potencia y la
          impedancia de entrada del transistor en sus tres configuraciones: emisor común, base
          común y colector común.
    \item Distinguir qué amplifica cada configuración y explicar por qué el emisor común
          pierde casi toda su ganancia con una carga pequeña y cómo un seguidor de emisor lo
          resuelve.
    \item Comparar una etapa de potencia push-pull en clase B y en clase AB: distorsión de
          cruce, potencia entregada a la carga y rendimiento.
\end{itemize}

%=======================================================================
\section{Contenido teórico}
%=======================================================================
\begin{itemize}
    \item El transistor como amplificador: punto Q, señal pequeña, capacitores de acoplo y de
          desacoplo.
    \item Ganancias de voltaje, corriente y potencia; impedancias de entrada y de salida;
          decibelios.
    \item Configuraciones en emisor común, base común y colector común; etapas de potencia
          clase B y clase AB.
\end{itemize}

\noindent\textbf{Duración:} dos sesiones. \textbf{Semana~1:} experimentos~1 y~2 (amplificadores
de voltaje). \textbf{Semana~2:} experimentos~3 y~4 (amplificadores de corriente y de
potencia).

%=======================================================================
\section{Equipos y materiales}
%=======================================================================
\noindent\textbf{Laboratorio:} Electrónica de Comunicaciones.

\vspace{0.25cm}
\noindent\textbf{Equipos y dispositivos:}
\begin{itemize}
    \item Tarjeta Analog Discovery Studio: riel fijo de \qty{+12}{\volt}, fuentes programables
          $V_+$ y $V_-$, generador de funciones \texttt{W1} y osciloscopio de dos canales.
\end{itemize}

\noindent\textbf{Software:}
\begin{itemize}
    \item Simulador Multisim o Proteus.
\end{itemize}

\vspace{0.25cm}
\needspace{6\baselineskip}
\noindent\textbf{Materiales y fungibles} (deben ser adquiridos por los alumnos):
\begin{itemize}
    % --- Semiconductores
    \item 02 transistores NPN BC548B, 03 NPN 2N2222A y 02 PNP 2N2907A (uno de repuesto de cada
          uno). Sus terminales están en distinto orden: revise la hoja de datos antes de armar.
    \item 03 diodos de señal 1N4148 (uno de repuesto).
    % --- Resistencias (la que mas disipa es RE = 1 k del exp. 3: 23 mW;
    %     RL = 47 Ohm del exp. 4 recibe ~60 mW)
    \item Resistencia a \qty{0.5}{\watt}: 01 de \qty{47}{\ohm}.
    \item Resistencias a \qty{0.25}{\watt}:
          02 de \qty{10}{\kilo\ohm},
          01 de \qty{1.8}{\kilo\ohm},
          01 de \qty{2.2}{\kilo\ohm},
          01 de \qty{390}{\ohm},
          04 de \qty{100}{\ohm},
          03 de \qty{1}{\kilo\ohm} y
          02 de \qty{47}{\kilo\ohm}.
    % --- Capacitores
    \item Capacitores electrolíticos de \qty{25}{\volt}: 03 de \qty{10}{\micro\farad},
          01 de \qty{47}{\micro\farad} y 03 de \qty{100}{\micro\farad}.
    % --- Montaje
    \item Protoboard.
    \item Jumpers macho-macho para protoboard.
    \item Alicates para cortar y pelar alambre.
    \item Multímetro.
    \item Lentes de protección y mandil.
\end{itemize}

\noindent Los experimentos~1 y~2 usan el circuito del punto medio del Laboratorio~5: si todavía
lo tiene armado, tráigalo. Guarde el seguidor del experimento~3 armado hasta terminar la
cascada.

\FloatBarrier
%=======================================================================
\section{Conocimientos previos}
%=======================================================================

%-----------------------------------------------------------------------
\subsection{¿Qué es amplificar?}
%-----------------------------------------------------------------------
En el Laboratorio~6 el transistor solo estaba en corte o en saturación. Para amplificar se lo
deja en la \textbf{región activa}, en un punto Q cerca del centro de la recta de carga
(Laboratorio~5), y se le suma una señal pequeña. Cada voltaje y cada corriente tiene entonces
dos partes: la \textbf{continua} (DC) del punto Q, que se escribe con mayúsculas ($V_C$,
$I_C$), y la \textbf{alterna} (AC) de la señal, con minúsculas ($v_c$, $i_c$):
\begin{equation*}
  v_C(t) = V_C + v_c(t).
\end{equation*}
Una pequeña variación de la corriente de base produce una variación $\beta$ veces mayor en la
corriente de colector. Esa es la base de los tres tipos de amplificación:
\begin{itemize}
  \item \textbf{De corriente:} $i_c = \beta\, i_b$; la carga recibe mucha más corriente de la
        que entrega la fuente de señal.
  \item \textbf{De voltaje:} la corriente $i_c$ pasa por una resistencia y, por la ley de Ohm,
        produce un voltaje mucho mayor que el de entrada.
  \item \textbf{De potencia:} como $p = v\,i$, si hay ganancia de voltaje, de corriente o de
        ambas, la carga recibe más potencia de señal que la que entrega la fuente. Esa potencia
        extra \textbf{no la crea el transistor}: sale de la fuente DC.
\end{itemize}

\noindent\textbf{Capacitores.} Los capacitores \textbf{de acoplo} ($C_{in}$ y $C_{out}$)
dejan pasar la señal y bloquean la continua, para que el generador y la carga no muevan el
punto Q. Los capacitores \textbf{de desacoplo} ($C_E$ y $C_B$) ponen un punto del circuito a
tierra solo para la señal (sección~\ref{sec:tierra}). A la frecuencia de trabajo
(\qty{1}{\kilo\hertz}) todos se comportan como cortocircuitos.

\noindent\textbf{Cómo se analiza un amplificador.} En dos pasos:
\begin{enumerate}[label=\arabic*), leftmargin=*, itemsep=4pt]
  \item \textbf{Análisis DC:} los capacitores son circuitos abiertos; se calcula el punto Q
        como en el Laboratorio~5 y, con él, la corriente de emisor $I_E$.
  \item \textbf{Análisis AC:} los capacitores y la fuente DC son cortocircuitos (la fuente
        $V_{CC}$ queda conectada a tierra). El transistor se reemplaza por su
        \textbf{resistencia dinámica de emisor}
        \begin{equation}
          r_e = \frac{\qty{26}{\milli\volt}}{I_E},
          \label{eq:re}
        \end{equation}
        que es la resistencia que la señal ``ve'' en el diodo base-emisor. Desde la base, una
        resistencia $R$ en el emisor se ve $\beta$ veces mayor: $\beta R$.
\end{enumerate}

%-----------------------------------------------------------------------
\subsection{Tierra para la señal}
\label{sec:tierra}
%-----------------------------------------------------------------------
Un terminal puede estar a varios voltios en DC y, aun así, estar \textbf{a tierra para la
señal}: basta con que en él el voltaje \emph{no varíe}. Hay dos formas de lograrlo.

\noindent\textbf{1. Con un capacitor de desacoplo.} La oposición de un capacitor al paso de la
corriente es su reactancia,
\begin{equation}
  X_C = \frac{1}{2\pi f C}.
  \label{eq:xc}
\end{equation}
En DC ($f = 0$) es infinita: el capacitor es un circuito abierto. A \qty{1}{\kilo\hertz}, un
capacitor de \qty{100}{\micro\farad} tiene $X_C \approx \qty{1.6}{\ohm}$: es casi un cable. Si
se pone en paralelo con una resistencia de \qty{390}{\ohm}, casi toda la corriente de la señal
se va por el capacitor y la resistencia deja de contar para la señal
(figura~\ref{fig:dcac}).

\begin{figure}[H]
  \centering
  \begin{subfigure}[b]{0.42\linewidth}
    \centering
    \begin{circuitikz}[scale=0.8, transform shape]
      \draw (0,0) node[npn](Q){};
      \coordinate (e2) at ($(Q.E)+(0,-1.7)$);
      \draw (Q.E) to[R, l_=$R_{E1}$, a^=\qty{100}{\ohm}] (e2)
                  to[R, l_=$R_{E2}$, a^=\qty{390}{\ohm}] ++(0,-1.7) node[ground]{};
      \draw[dashed, black!50] (e2) -- ++(1.4,0) coordinate(ce);
      \draw[black!50] (ce) to[C, color=black!50] ++(0,-1.7) -- ++(0,-0.01);
      \node[black!60, right, align=left, font=\small] at ($(ce)+(0.3,-0.85)$) {$C_E$\\abierto};
      \draw (Q.E) ++(-0.2,-0.2) node[left, font=\small]{emisor};
    \end{circuitikz}
    \caption{En DC: el emisor ve $R_{E1} + R_{E2} = \qty{490}{\ohm}$. El punto Q queda
             estable.}
  \end{subfigure}
  \hfill
  \begin{subfigure}[b]{0.42\linewidth}
    \centering
    \begin{circuitikz}[scale=0.8, transform shape]
      \draw (0,0) node[npn](Q){};
      \coordinate (e2) at ($(Q.E)+(0,-1.7)$);
      \draw (Q.E) to[R, l_=$R_{E1}$, a^=\qty{100}{\ohm}] (e2);
      \draw[black!40] (e2) to[R, color=black!40] ++(0,-1.7) coordinate(g);
      \draw (e2) to[short, *-] ++(1.4,0) -- ++(0,-1.7) -- (g) node[ground]{};
      \node[right, align=left, font=\small] at ($(e2)+(1.5,-0.85)$) {$C_E$ = cable\\[1pt]($R_{E2}$ no\\cuenta)};
      \draw (Q.E) ++(-0.2,-0.2) node[left, font=\small]{emisor};
    \end{circuitikz}
    \caption{En AC (\qty{1}{\kilo\hertz}): la señal solo ve $R_{E1} = \qty{100}{\ohm}$.}
  \end{subfigure}
  \caption{El emisor del amplificador del experimento~1 visto en DC y en AC. Con
           \qty{490}{\ohm} el punto Q casi no depende de $\beta$; con \qty{100}{\ohm} la
           ganancia es mayor.}
  \label{fig:dcac}
\end{figure}

\noindent Así se consiguen las dos cosas a la vez: un punto Q estable, porque en DC hay mucha
resistencia de emisor, y una buena ganancia, porque la señal ve poca. ¿Por qué no desacoplar
toda la resistencia de emisor? Porque entonces la ganancia dependería solo de $r_e$, que cambia
con la temperatura, de un transistor a otro e incluso durante cada ciclo de la señal: la
ganancia sería inestable y la salida saldría deformada. Dejar $R_{E1}$ sin desacoplar es un
término medio.

\noindent\textbf{2. Con la fuente de alimentación.} Una fuente DC ideal mantiene su voltaje
fijo: en ella la señal vale cero. Por eso $V_{CC}$, aunque esté a \qty{12}{\volt}, es tierra
para la señal, y en el análisis AC se reemplaza por un cable a tierra.

\noindent\textbf{El nombre de cada configuración.} El terminal \textbf{común} es el que
comparten la entrada y la salida: el que queda a tierra para la señal.

\begin{table}[H]
  \centering
  \renewcommand{\arraystretch}{1.3}
  \begin{tabularx}{\linewidth}{|>{\raggedright\arraybackslash}m{3.0cm}|*{4}{>{\centering\arraybackslash}X|}}
    \hline
    \textbf{Configuración} & \textbf{Entra por} & \textbf{Sale por}
      & \textbf{Terminal común} & \textbf{Lo pone a tierra para la señal} \\ \hline
    Emisor común   & base   & colector & emisor   & $C_E$ \\ \hline
    Base común     & emisor & colector & base     & $C_B$ \\ \hline
    Colector común & base   & emisor   & colector & la fuente $V_{CC}$ \\ \hline
  \end{tabularx}
  \caption{Las tres configuraciones del transistor bipolar.}
  \label{tab:configuraciones}
\end{table}

%-----------------------------------------------------------------------
\subsection{Las ganancias y las impedancias}
\label{sec:ganancias}
%-----------------------------------------------------------------------
Un amplificador se describe con cinco números (figura~\ref{fig:bloque}):

\begin{figure}[H]
  \centering
  \begin{circuitikz}[scale=0.85, transform shape]
    \draw (0,0) to[sV, l=$v_s$] (0,2.4) to[R, l=$R_s$, -o] (2.4,2.4)
          node[above]{$v_{in}$};
    \draw (2.4,2.4) to[short, i_=$i_{in}$] (3.6,2.4) to[R, l_=$Z_{in}$] (3.6,0);
    \draw[thick, fill=black!5] (4.4,-0.4) rectangle (7.0,2.8);
    \node[align=center] at (5.7,1.2) {Amplificador\\[2pt] $A_v,\ A_i,\ A_p$};
    \draw (3.6,2.4) -- (4.4,2.4);
    \draw (7.0,2.4) to[R, l=$Z_{out}$] (8.6,2.4) to[short, -o, i=$i_o$] (9.6,2.4)
          node[above]{$v_o$} to[R, l=$R_L$] (9.6,0);
    \draw (0,0) -- (9.6,0);
    \draw (0,0) node[ground]{};
  \end{circuitikz}
  \caption{Un amplificador visto desde afuera. La fuente ve la impedancia de entrada
           $Z_{in}$; la carga ve una fuente con impedancia de salida $Z_{out}$.}
  \label{fig:bloque}
\end{figure}

\begin{table}[H]
  \centering
  \renewcommand{\arraystretch}{2.0}
  \begin{tabularx}{\linewidth}{|>{\raggedright\arraybackslash}m{4.2cm}|>{\centering\arraybackslash}m{3.6cm}|>{\raggedright\arraybackslash}X|}
    \hline
    \textbf{Magnitud} & \textbf{Definición} & \textbf{Qué indica} \\ \hline
    Ganancia de voltaje & $A_v = \dfrac{v_o}{v_{in}}$ & Cuántas veces crece el voltaje. Si es
      negativa, la salida está invertida (desfasada \ang{180}). \\ \hline
    Ganancia de corriente & $A_i = \dfrac{i_o}{i_{in}}$ & Cuántas veces crece la corriente;
      $i_o = v_o/R_L$ e $i_{in} = v_{in}/Z_{in}$. \\ \hline
    Ganancia de potencia & $A_p = \dfrac{p_o}{p_{in}} = |A_v|\,A_i$ & Cuántas veces crece la
      potencia de la señal. \\ \hline
    Impedancia de entrada & $Z_{in} = \dfrac{v_{in}}{i_{in}}$ & Cuánto carga el amplificador a
      la fuente de señal. Conviene que sea alta. \\ \hline
    Impedancia de salida & $Z_{out}$ & Cuánto cae la salida al conectar la carga. Conviene que
      sea baja. \\ \hline
  \end{tabularx}
  \caption{Las magnitudes de un amplificador.}
  \label{tab:magnitudes}
\end{table}

\noindent Las tres ganancias están relacionadas: basta con medir $A_v$ y $Z_{in}$.
\begin{equation}
  A_i = \frac{v_o/R_L}{v_{in}/Z_{in}} = |A_v|\,\frac{Z_{in}}{R_L}
  \qquad\qquad
  A_p = |A_v|\,A_i = A_v^2\,\frac{Z_{in}}{R_L}
  \label{eq:aiap}
\end{equation}
\noindent La carga y la impedancia de salida forman un divisor de tensión, así que la ganancia
con carga es menor que sin ella:
\begin{equation}
  A_{v(\text{con } R_L)} = A_{v(\text{sin carga})}\,\frac{R_L}{R_L + Z_{out}}.
  \label{eq:divisor}
\end{equation}

\noindent\textbf{En decibelios.} $A_v(\unit{dB}) = 20\log_{10}|A_v|$ y
$A_p(\unit{dB}) = 10\log_{10} A_p$. Una ganancia de potencia de 100 son \qty{20}{dB}.

\noindent\textbf{Cómo se mide $\boldsymbol{Z_{in}}$} (en el simulador y en la tarjeta). Se pone
una resistencia conocida $R_s$ entre el generador y la entrada (figura~\ref{fig:bloque}) y se
miden los dos voltajes pico $v_s$ y $v_{in}$. Por $R_s$ pasa $i_{in} = (v_s - v_{in})/R_s$,
de modo que
\begin{equation}
  Z_{in} = R_s\,\frac{v_{in}}{v_s - v_{in}}.
  \label{eq:zin}
\end{equation}
Funciona mejor si $R_s$ es del orden de $Z_{in}$: si es mucho menor, $v_s$ y $v_{in}$ casi no
se diferencian.

%-----------------------------------------------------------------------
\subsection{Amplificador de voltaje: emisor común}
%-----------------------------------------------------------------------
La señal entra por la base y sale por el colector; el emisor va a tierra para la señal a
través de $C_E$ (figura~\ref{fig:ce-teo}). Solo la parte $R_{E1}$ de la resistencia de emisor
queda sin desacoplar: baja la ganancia, pero la hace estable y menos dependiente de $r_e$.

\par\needspace{8cm}\vspace{8pt}\noindent
\begin{minipage}[c]{0.56\linewidth}
  \centering
  \begin{circuitikz}[scale=0.8, transform shape]
    \draw (4.6,0) node[npn](Q){};
    \coordinate (gnd) at (0,-4.2);
    \coordinate (top) at (0,3.4);
    \coordinate (b) at ($(Q.B)+(-0.9,0)$);
    \draw (b |- top) to[R, l_=$R_1$] (b) to[R, l_=$R_2$] (b |- gnd);
    \draw (b) to[short, *-] (Q.B);
    \coordinate (in) at ($(b)+(-1.8,0)$);
    \draw (in) to[C, l=$C_{in}$] (b);
    \draw (in |- gnd) to[sV, l=$v_{in}$] (in);
    \coordinate (c) at ($(Q.C)+(0,0.5)$);
    \draw (c |- top) to[R, l_=$R_C$] (c) -- (Q.C);
    \draw (b |- top) -- (c |- top) node[circ]{} -- ++(0.8,0) node[right]{$V_{CC}$};
    \coordinate (e2) at ($(Q.E)+(0,-1.7)$);
    \draw (Q.E) to[R, l_=$R_{E1}$] (e2) to[R, l_=$R_{E2}$] (e2 |- gnd);
    \draw (e2) to[short, *-] ++(1.3,0) coordinate(ce) to[C, l=$C_E$] (ce |- gnd);
    \draw (c) to[C, l=$C_{out}$, *-] ++(2.4,0) coordinate(o)
              to[R, l=$R_L$, *-] (o |- gnd);
    \draw (o) to[short, -o] ++(0.6,0) node[right]{$v_o$};
    \draw (in |- gnd) -- (o |- gnd);
    \draw (Q.E |- gnd) node[ground]{};
  \end{circuitikz}
  \captionof{figure}{Emisor común: la señal entra por la base y sale por el colector.}
  \label{fig:ce-teo}
\end{minipage}\hfill
\begin{minipage}[c]{0.42\linewidth}
  \begin{align}
    A_v &= -\frac{R_C \pll R_L}{r_e + R_{E1}} \label{eq:avce}\\[8pt]
    Z_{in} &= R_1 \pll R_2 \pll \beta\,(r_e + R_{E1}) \label{eq:zince}\\[8pt]
    Z_{out} &\approx R_C \label{eq:zoutce}
  \end{align}
\end{minipage}
\par\vspace{10pt}\noindent
donde $R_a \pll R_b = R_a R_b/(R_a + R_b)$. El signo menos indica que la salida está
invertida. Es la configuración más usada porque da ganancia de voltaje \emph{y} de corriente,
pero su $Z_{out}$ es alta: con una carga pequeña, $R_C \pll R_L$ se hace pequeño y la ganancia
se desploma.

\needspace{14\baselineskip}
\noindent\textbf{Ejemplo resuelto} (valores distintos a los de la práctica). Un emisor común
con $I_E = \qty{1.3}{\milli\ampere}$, $R_C = \qty{3.3}{\kilo\ohm}$,
$R_{E1} = \qty{150}{\ohm}$, $R_1 \pll R_2 = \qty{4}{\kilo\ohm}$, $\beta = 150$ y
$R_L = \qty{4.7}{\kilo\ohm}$.
\begin{table}[H]
  \centering
  \renewcommand{\arraystretch}{2.1}
  \begin{tabularx}{\linewidth}{|>{\centering\arraybackslash}m{2.4cm}|>{\centering\arraybackslash}X|}
    \hline
    \textbf{Paso} & \textbf{Cálculo} \\ \hline
    $r_e$ & $\dfrac{\qty{26}{\milli\volt}}{\qty{1.3}{\milli\ampere}} = \qty{20}{\ohm}$ \\ \hline
    $A_v$ & $R_C \pll R_L = \qty{1.94}{\kilo\ohm}$, \qquad
            $A_v = -\dfrac{1940}{20 + 150} = \num{-11.4}$ \quad (\qty{21.1}{dB}) \\ \hline
    $Z_{in}$ & $\beta\,(r_e + R_{E1}) = 150 \times 170 = \qty{25.5}{\kilo\ohm}$, \qquad
               $Z_{in} = \num{4} \pll \qty{25.5}{\kilo\ohm} = \qty{3.46}{\kilo\ohm}$ \\ \hline
    $A_i$, $A_p$ & $A_i = \num{11.4} \times \dfrac{\num{3.46}}{\num{4.7}} = \num{8.4}$, \qquad
                   $A_p = \num{11.4} \times \num{8.4} = 96$ \quad (\qty{19.8}{dB}) \\ \hline
  \end{tabularx}
  \caption{Ejemplo de cálculo de un emisor común.}
  \label{tab:ejemplo}
\end{table}

%-----------------------------------------------------------------------
\subsection{Amplificador de voltaje: base común}
%-----------------------------------------------------------------------
La señal entra por el \textbf{emisor} y sale por el colector; la base va a tierra para la
señal a través de $C_B$ (figura~\ref{fig:cb-teo}). La polarización es la misma del emisor
común: solo cambia por dónde entra la señal. Como la señal mueve directamente la corriente de
emisor, la corriente de entrada es casi toda $i_e$ y la de salida es $i_c = \alpha\,i_e$, con
$\alpha = \beta/(\beta + 1) \approx 1$.

\par\needspace{7.5cm}\vspace{8pt}\noindent
\begin{minipage}[c]{0.56\linewidth}
  \centering
  \begin{circuitikz}[scale=0.8, transform shape]
    \draw (3.4,0) node[npn](Q){};
    \coordinate (gnd) at (0,-4.0);
    \coordinate (top) at (0,3.0);
    \coordinate (b) at ($(Q.B)+(-0.9,0)$);
    \draw (b |- top) to[R, l_=$R_1$] (b) to[R, l_=$R_2$] (b |- gnd);
    \draw (b) to[short, *-] (Q.B);
    \draw (b) -- ++(-1.3,0) coordinate(cb) to[C, l_=$C_B$] (cb |- gnd);
    \coordinate (c) at ($(Q.C)+(0,0.5)$);
    \draw (c |- top) to[R, l_=$R_C$] (c) -- (Q.C);
    \draw (cb |- top) -- (c |- top) node[circ]{} -- ++(0.8,0) node[right]{$V_{CC}$};
    \coordinate (e) at ($(Q.E)+(0,-1.1)$);
    \draw (Q.E) -- (e) to[R, l_=$R_E$] (e |- gnd);
    \draw (e) to[C, l_=$C_{in}$, *-] ++(1.7,0) coordinate(in)
              to[R, l_=$R_s$] ++(1.7,0) coordinate(s)
              to[sV, l_=$v_s$] (s |- gnd);
    \node[above] at (e) {$v_{in}$};
    \draw (c) to[C, l=$C_{out}$, *-] ++(1.8,0) -- ++(2.6,0) coordinate(o)
              to[short, *-o] ++(0.7,0) node[right]{$v_o$};
    \draw (o) to[R, l=$R_L$] (o |- gnd);
    \draw (cb |- gnd) -- (o |- gnd);
    \draw (b |- gnd) node[ground]{};
  \end{circuitikz}
  \captionof{figure}{Base común: la señal entra por el emisor y sale por el colector.}
  \label{fig:cb-teo}
\end{minipage}\hfill
\begin{minipage}[c]{0.42\linewidth}
  \begin{align}
    A_v &= \frac{R_C \pll R_L}{r_e} \label{eq:avcb}\\[8pt]
    Z_{in} &= R_E \pll r_e \approx r_e \label{eq:zincb}\\[8pt]
    A_i &\approx \alpha\,\frac{R_C}{R_C + R_L} < 1 \label{eq:aicb}\\[8pt]
    Z_{out} &\approx R_C \label{eq:zoutcb}
  \end{align}
\end{minipage}
\par\vspace{10pt}\noindent
La ganancia de voltaje es alta y \textbf{sin inversión}, pero la impedancia de entrada es de
apenas unos ohmios. Una fuente de señal con resistencia interna $R_s$ pierde casi todo su voltaje
en ella: de la fuente a la salida, la ganancia queda en
\begin{equation}
  \frac{v_o}{v_s} \approx \frac{R_C \pll R_L}{R_s + r_e}.
  \label{eq:avscb}
\end{equation}
Por eso la base común se usa con fuentes de impedancia muy baja y en altas frecuencias (radio,
amplificador cascodo), donde su mejor respuesta en frecuencia compensa esa desventaja.

%-----------------------------------------------------------------------
\subsection{Amplificador de corriente: colector común (seguidor de emisor)}
%-----------------------------------------------------------------------
La señal entra por la base y sale por el \textbf{emisor}; el colector va directo a $V_{CC}$
(figura~\ref{fig:cc-teo}). El emisor ``sigue'' a la base a \qty{0.7}{\volt} por debajo, así que
$A_v$ es casi 1 y sin inversión. Su utilidad está en las impedancias: una carga $R_L$ en el
emisor se ve desde la base $\beta$ veces mayor, y la salida tiene impedancia muy baja.

\par\vspace{8pt}\noindent
\begin{minipage}[c]{0.56\linewidth}
  \centering
  \begin{circuitikz}[scale=0.8, transform shape]
    \draw (5.4,0) node[npn](Q){};
    \coordinate (gnd) at (0,-3.6);
    \coordinate (top) at (0,2.8);
    \coordinate (b) at ($(Q.B)+(-0.9,0)$);
    \draw (b |- top) to[R, l_=$R_1$] (b) to[R, l_=$R_2$] (b |- gnd);
    \draw (b) to[short, *-] (Q.B);
    \coordinate (in) at ($(b)+(-1.8,0)$);
    \coordinate (s)  at ($(in)+(-1.6,0)$);
    \draw (in) to[C, l_=$C_{in}$] (b);
    \draw (s) to[R, l=$R_s$] (in);
    \draw (s |- gnd) to[sV, l=$v_s$] (s);
    \draw (Q.C) -- (Q.C |- top);
    \draw (b |- top) -- (Q.C |- top) node[circ]{} -- ++(0.8,0) node[right]{$V_{CC}$};
    \coordinate (e) at ($(Q.E)+(0,-0.5)$);
    \draw (Q.E) -- (e) to[R, l_=$R_E$] (e |- gnd);
    \draw (e) to[C, l=$C_{out}$, *-] ++(2.4,0) coordinate(o)
              to[R, l=$R_L$, *-] (o |- gnd);
    \draw (o) to[short, -o] ++(0.6,0) node[right]{$v_o$};
    \draw (s |- gnd) -- (o |- gnd);
    \draw (Q.E |- gnd) node[ground]{};
  \end{circuitikz}
  \captionof{figure}{Colector común: la señal entra por la base y sale por el emisor.}
  \label{fig:cc-teo}
\end{minipage}\hfill
\begin{minipage}[c]{0.42\linewidth}
  \begin{align}
    A_v &= \frac{R_E \pll R_L}{r_e + R_E \pll R_L} \approx 1 \label{eq:avcc}\\[8pt]
    Z_{in} &= R_1 \pll R_2 \pll \beta\,(r_e + R_E \pll R_L) \label{eq:zincc}\\[8pt]
    Z_{out} &\approx r_e + \frac{R_s \pll R_1 \pll R_2}{\beta} \label{eq:zoutcc}
  \end{align}
\end{minipage}
\par\vspace{10pt}\noindent
Con $Z_{in}$ grande y $R_L$ pequeña, por~\eqref{eq:aiap} la ganancia de corriente es grande
aunque $A_v < 1$. Por eso se usa como \textbf{etapa separadora} (\emph{buffer}) entre un
emisor común y una carga pequeña.

\noindent\textbf{Límite de la salida.} En el semiciclo negativo el transistor solo puede
apagarse; quien empuja corriente por la carga es $I_E$. La salida se recorta en
\begin{equation}
  v_{o(\min)} = -I_E\,(R_E \pll R_L).
  \label{eq:recorte}
\end{equation}

%-----------------------------------------------------------------------
\subsection{Amplificador de potencia: push-pull clase B y clase AB}
%-----------------------------------------------------------------------
Para entregar potencia a una carga pequeña (un altavoz, un motor) se usan dos seguidores de
emisor complementarios, un NPN y un PNP, con fuente doble $\pm V_{CC}$
(figura~\ref{fig:pp-teo}). El NPN conduce en el semiciclo positivo y el PNP en el negativo:
cada uno ``empuja'' (\emph{push}) o ``jala'' (\emph{pull}) la corriente de la carga.

\begin{figure}[!tbp]
  \centering
  \begin{subfigure}[b]{0.40\linewidth}
    \centering
    \begin{circuitikz}[scale=0.68, transform shape]
      \draw (5,1.8) node[npn](Q1){};
      \draw (5,-1.8) node[pnp](Q2){};
      \draw (Q1.E) -- (Q2.E);
      \coordinate (o) at ($(Q1.E)!0.5!(Q2.E)$);
      \coordinate (bb) at ($(Q1.B)!0.5!(Q2.B)+(-0.8,0)$);
      \draw (Q1.B) -| (bb) |- (Q2.B);
      \draw (Q1.C) -- ++(0,0.5) node[above]{$+V_{CC}$};
      \draw (Q2.C) -- ++(0,-0.5) node[below]{$-V_{CC}$};
      \draw (o) to[short, *-] ++(1.6,0) coordinate(ro)
                to[R, l=$R_L$] ++(0,-2.2) node[ground]{};
      \draw (ro) to[short, *-o] ++(0.6,0) node[above]{$v_o$};
      \draw (bb) to[short, *-] ++(-0.9,0) coordinate(s)
                 to[sV, l_=$v_{in}$] ++(0,-2.2) node[ground]{};
      \node[right, font=\small] at (Q1.east) {NPN};
      \node[right, font=\small] at (Q2.east) {PNP};
    \end{circuitikz}
    \caption{Clase B.}
  \end{subfigure}
  \hfill
  \begin{subfigure}[b]{0.52\linewidth}
    \centering
    \begin{circuitikz}[scale=0.68, transform shape]
      \draw (5,1.8) node[npn](Q1){};
      \draw (5,-1.8) node[pnp](Q2){};
      \draw (Q1.E) -- (Q2.E);
      \coordinate (o) at ($(Q1.E)!0.5!(Q2.E)$);
      \coordinate (b1) at ($(Q1.B)+(-0.9,0)$);
      \coordinate (b2) at ($(Q2.B)+(-0.9,0)$);
      \coordinate (m)  at ($(b1)!0.5!(b2)$);
      \draw (Q1.B) -- (b1);
      \draw (Q2.B) -- (b2);
      \draw (b1) to[D, l=$D_1$, *-*] (m) to[D, l=$D_2$, -*] (b2);
      \draw (Q1.C) |- ($(b1)+(0,2.0)$) coordinate(tp);
      \draw (tp) to[R, l_=$R_{B1}$] (b1);
      \draw (Q2.C) |- ($(b2)+(0,-2.0)$) coordinate(bt);
      \draw (b2) to[R, l_=$R_{B2}$] (bt);
      \node[above] at (Q1.C |- tp) {$+V_{CC}$};
      \node[below] at (Q2.C |- bt) {$-V_{CC}$};
      \draw (o) to[short, *-] ++(1.6,0) coordinate(ro)
                to[R, l=$R_L$] ++(0,-2.2) node[ground]{};
      \draw (ro) to[short, *-o] ++(0.6,0) node[above]{$v_o$};
      \coordinate (in) at ($(m)+(-1.9,0)$);
      \draw (in) to[C, l=$C_{in}$] (m);
      \draw (in) to[sV, l_=$v_{in}$] ++(0,-2.2) node[ground]{};
    \end{circuitikz}
    \caption{Clase AB.}
  \end{subfigure}
  \caption{Etapa push-pull complementaria. El NPN conduce en el semiciclo positivo y el PNP en
           el negativo; en la clase AB los diodos polarizan a los dos al borde de conducir.}
  \label{fig:pp-teo}
\end{figure}

\noindent\textbf{Clase B.} Las bases van juntas a la entrada. Sin señal, ninguno conduce: no
hay consumo en reposo, pero mientras $|v_{in}| < \qty{0.7}{\volt}$ los dos están cortados y la
salida se queda en cero. Esa ``muesca'' en cada cruce por cero es la \textbf{distorsión de
cruce}.

\noindent\textbf{Clase AB.} Dos diodos entre las bases (figura~\ref{fig:pp-teo}b) les dan
\qty{0.7}{\volt} a cada una, de modo que los dos transistores quedan al borde de conducir. La
distorsión de cruce desaparece a cambio de una pequeña corriente de reposo.

\noindent\textbf{Potencias.} Para una salida senoidal de pico $V_p$ en la carga $R_L$:
\begin{align}
  P_{out} &= \frac{V_p^2}{2R_L} = \frac{V_{o(\text{rms})}^2}{R_L} \label{eq:pout}\\[4pt]
  P_{DC} &= V_{CC}\,I_{+} + |V_{EE}|\,I_{-} \;\approx\; \frac{2\,V_{CC}\,V_p}{\pi\,R_L}
     \quad\text{(clase B)} \label{eq:pdc}\\[4pt]
  \eta &= \frac{P_{out}}{P_{DC}} = \frac{\pi}{4}\,\frac{V_p}{V_{CC}} \quad\text{(clase B)}
     \label{eq:eta}
\end{align}
donde $I_+$ e $I_-$ son las corrientes \textbf{promedio} (las que lee un multímetro en DC) de
cada fuente. El rendimiento máximo de la clase B es $\pi/4 = \qty{78.5}{\percent}$, cuando
$V_p = V_{CC}$. En la clase AB hay que sumarle a $P_{DC}$ el consumo de los diodos y sus
resistencias. Cada transistor disipa como máximo
$P_{D(\max)} = V_{CC}^2/(\pi^2 R_L)$.

\needspace{10\baselineskip}
\noindent\textbf{Resumen de las configuraciones.}
\begin{table}[H]
  \centering
  \renewcommand{\arraystretch}{1.3}
  \begin{tabularx}{\linewidth}{|>{\raggedright\arraybackslash}m{2.6cm}|*{4}{>{\centering\arraybackslash}X|}}
    \hline
    & \textbf{Emisor común} & \textbf{Base común} & \textbf{Colector común} & \textbf{Push-pull} \\ \hline
    Amplifica sobre todo & voltaje & voltaje & corriente & potencia \\ \hline
    $A_v$ & alta, invertida & alta, sin invertir & $\approx 1$ & $\approx 1$ \\ \hline
    $A_i$ & media & $< 1$ & alta & alta \\ \hline
    $Z_{in}$ & media (\unit{\kilo\ohm}) & muy baja (\unit{\ohm}) & alta (decenas de \unit{\kilo\ohm}) & media \\ \hline
    $Z_{out}$ & alta ($\approx R_C$) & alta ($\approx R_C$) & muy baja (\unit{\ohm}) & muy baja \\ \hline
    Uso típico & etapa de entrada & altas frecuencias & separador (\emph{buffer}) & etapa de salida \\ \hline
    Experimento & 1 & 2 & 3 & 4 \\ \hline
  \end{tabularx}
  \caption{Comparación de las configuraciones de la práctica.}
  \label{tab:comparacion}
\end{table}

\FloatBarrier
\needspace{14\baselineskip}
%=======================================================================
\section{Desarrollo del laboratorio}
%=======================================================================
La práctica dura dos semanas. Cada experimento tiene los mismos cuatro pasos. Los pasos~1 y~2
(\textbf{Datos} y \textbf{Calcule}) se traen resueltos de casa, junto con la simulación en
Multisim o Proteus; los pasos~3 y~4 (\textbf{Arme y mida} y \textbf{Anote}) se hacen en el
laboratorio. Use $\beta = 200$, $V_{BE} = \qty{0.7}{\volt}$ y
$r_e = \qty{26}{\milli\volt}/I_E$.

\begin{table}[H]
  \centering
  \renewcommand{\arraystretch}{1.35}
  \begin{tabularx}{\linewidth}{|>{\centering\arraybackslash}m{1.1cm}|>{\centering\arraybackslash}m{0.8cm}|>{\raggedright\arraybackslash}m{2.7cm}|>{\raggedright\arraybackslash}X|>{\raggedright\arraybackslash}X|>{\centering\arraybackslash}m{1.3cm}|}
    \hline
    \textbf{Sem.} & \textbf{Exp.} & \textbf{Circuito} & \textbf{En casa: calcule}
      & \textbf{En el laboratorio: mida} & \textbf{Fuente} \\ \hline
    \multirow{2}{*}{1} & 1 & Emisor común (voltaje)
      & Punto Q, $r_e$, $A_v$, $Z_{in}$, $A_i$ y $A_p$ con $R_L = \qty{10}{\kilo\ohm}$ y
        \qty{100}{\ohm}
      & $V_B$, $V_E$, $V_C$; $v_s$, $v_{in}$ y $v_o$ con cada carga
      & \qty{+12}{\volt} \\ \cline{2-6}
    & 2 & Base común (voltaje)
      & $A_v$, $Z_{in}$, $A_i$ y $A_p$; ganancia desde la fuente
      & $v_s$, $v_{in}$ y $v_o$ sin carga y con $R_L = \qty{10}{\kilo\ohm}$
      & \qty{+12}{\volt} \\ \hline
    \multirow{2}{*}{2} & 3 & Colector común (corriente) y cascada
      & Punto Q, $A_v$, $Z_{in}$, $A_i$, $A_p$ y el límite de recorte; $A_v$ de la cascada
      & La carga sin y con transistor; el recorte; $v_o$ de la cascada en \qty{100}{\ohm}
      & \qty{+12}{\volt} \\ \cline{2-6}
    & 4 & Push-pull clase B y AB (potencia)
      & $P_{out}$, $P_{DC}$, $\eta$ y $P_{D(\max)}$
      & Formas de onda; $v_o$, $I_+$ e $I_-$; $A_p$ en clase AB
      & $\pm\qty{5}{\volt}$ \\ \hline
  \end{tabularx}
  \caption{Resumen de la práctica.}
  \label{tab:resumen}
\end{table}

\noindent\textbf{Cómo simular las ganancias.} En Multisim o Proteus ponga una fuente senoidal
de \qty{1}{\kilo\hertz} con la amplitud indicada en cada experimento y haga un
\textbf{análisis transitorio} de \qty{0}{} a \qty{20}{\milli\second} (deje pasar los primeros
milisegundos, mientras se cargan los capacitores). Mida con los cursores, o con un
osciloscopio virtual, los valores pico de $v_s$, $v_{in}$ y $v_o$, y calcule $A_v$, $Z_{in}$
con~\eqref{eq:zin}, $A_i$ y $A_p$ con~\eqref{eq:aiap}. Los voltajes DC del punto Q se obtienen
con un análisis de punto de operación (\emph{DC operating point}).

\noindent\textbf{Cómo medir con la tarjeta.} El osciloscopio tiene dos canales, así que cada
medición se hace en dos capturas, siempre con acoplamiento DC y los canales \texttt{1$-$} y
\texttt{2$-$} a \texttt{GND}:
\begin{itemize}
  \item \textbf{Captura A:} CH1 en $v_s$ (la salida de \texttt{W1}) y CH2 en $v_{in}$ (después de
        $R_s$). Da $Z_{in}$ con~\eqref{eq:zin}.
  \item \textbf{Captura B:} CH1 en $v_{in}$ y CH2 en $v_o$. Da $A_v$.
\end{itemize}
Use la herramienta \emph{Measurements} de WaveForms para leer la amplitud o el valor pico a
pico de cada canal; promedie varias adquisiciones (\emph{Average}) cuando la señal sea de pocos
milivoltios.

%-----------------------------------------------------------------------
\subsection{Experimento 1: amplificador de voltaje en emisor común}
%-----------------------------------------------------------------------
Es el circuito del punto medio del Laboratorio~5, ahora con una resistencia $R_s$ para medir
$Z_{in}$ y un capacitor $C_{out}$ que lleva la señal a la carga.

\paso{1}{Datos}
$V_{CC} = \qty{12}{\volt}$ (riel fijo); \texttt{W1}: senoidal de \qty{1}{\kilo\hertz} y
\qty{50}{\milli\volt} de amplitud (pico); transistor BC548B. Los valores de los componentes
están en la figura~\ref{fig:ce}. Se prueban tres cargas: sin $R_L$,
$R_L = \qty{10}{\kilo\ohm}$ (la entrada de otra etapa) y $R_L = \qty{100}{\ohm}$ (una carga
pequeña).

\begin{figure}[H]
  \centering
  \begin{circuitikz}[scale=0.85, transform shape]
    \draw (7.4,0) node[npn](Q){};
    \node[right] at (Q.east) {BC548B};
    \coordinate (gnd) at (0,-4.4);
    \coordinate (top) at (0,3.8);
    \coordinate (b) at ($(Q.B)+(-1.0,0)$);
    % divisor de base
    \draw (b |- top) to[R, l_=$R_1$, a^=\qty{10}{\kilo\ohm}] (b);
    \draw (b) to[R, l_=$R_2$, a^=\qty{1.8}{\kilo\ohm}] (b |- gnd);
    \draw (b) to[short, *-] (Q.B);
    % entrada
    \coordinate (in) at ($(b)+(-2.4,0)$);
    \coordinate (s)  at ($(in)+(-2.6,0)$);
    \draw (in) to[C, l=$C_{in}$, a=\qty{10}{\micro\farad}] (b);
    \draw (s) to[R, l=$R_s$, a=\qty{1}{\kilo\ohm}, -*] (in);
    \node[above] at (in) {$v_{in}$};
    \draw (s |- gnd) to[sV, l=\texttt{W1}] (s);
    \node[above] at (s) {$v_s$};
    % colector
    \coordinate (c) at ($(Q.C)+(0,0.6)$);
    \draw (c |- top) to[R, l_=$R_C$, a^=\qty{2.2}{\kilo\ohm}] (c) -- (Q.C);
    \draw (b |- top) -- (c |- top) node[circ]{} -- ++(1.8,0) node[right]{$V_{CC} = \qty{+12}{\volt}$};
    % emisor
    \coordinate (e2) at ($(Q.E)+(0,-1.8)$);
    \draw (Q.E) to[R, l_=$R_{E1}$, a^=\qty{100}{\ohm}] (e2)
                to[R, l_=$R_{E2}$, a^=\qty{390}{\ohm}] (e2 |- gnd);
    \draw (e2) to[short, *-] ++(2.3,0) coordinate(ce)
               to[C, l={$C_E$ \qty{100}{\micro\farad}}] (ce |- gnd);
    % salida
    \draw (c) to[C, l=$C_{out}$, a=\qty{10}{\micro\farad}, *-] ++(3.6,0) coordinate(o)
              to[R, l=$R_L$, *-] (o |- gnd);
    \draw (o) to[short, -o] ++(1.0,0) node[right]{$v_o$};
    % tierra
    \draw (s |- gnd) -- (o |- gnd);
    \draw (Q.E |- gnd) node[ground]{};
  \end{circuitikz}
  \caption{Experimento 1: amplificador en emisor común. $R_s$ solo sirve para medir $Z_{in}$.
           La franja de $C_{out}$ va hacia la carga (el colector está a unos
           \qty{7}{\volt}) y la de $C_E$ hacia tierra.}
  \label{fig:ce}
\end{figure}

\paso{2}{Calcule (en casa)}
\begin{enumerate}[label=\alph*), leftmargin=*, itemsep=3pt]
  \item El punto Q: $V_{TH}$, $R_{TH}$, $I_E$, $V_B$, $V_E$ y $V_C$ (como en el
        Laboratorio~5). Con $I_E$, calcule $r_e$ con~\eqref{eq:re}.
  \item $A_v$ sin carga y con cada $R_L$ con~\eqref{eq:avce} (sin carga,
        $R_C \pll R_L = R_C$), y $Z_{in}$ con~\eqref{eq:zince}.
  \item $A_i$ y $A_p$ con~\eqref{eq:aiap} para $R_L = \qty{10}{\kilo\ohm}$ y
        $R_L = \qty{100}{\ohm}$; exprese $A_v$ y $A_p$ en decibelios. Siga el ejemplo de la
        tabla~\ref{tab:ejemplo}.
  \item Simule el circuito con las tres cargas y anote $v_s$, $v_{in}$ y $v_o$.
\end{enumerate}

\paso{3}{Arme y mida (en el laboratorio)}
\begin{enumerate}[label=\alph*), ref=\alph*, leftmargin=*, itemsep=3pt]
  \item Arme el circuito sin $R_L$ y aliméntelo con el riel de \qty{+12}{\volt}. Con
        \texttt{W1} apagado, mida $V_B$, $V_E$ y $V_C$ con el multímetro. Si $V_C$ no está
        entre \qtyrange{6}{8}{\volt}, revise el montaje antes de seguir.
  \item Encienda \texttt{W1} y haga las capturas A y B. Compruebe en la captura B que $v_o$
        está invertida respecto de $v_{in}$ y que no está recortada.
  \item \label{paso:carga}Repita la captura B con $R_L = \qty{10}{\kilo\ohm}$ y con
        $R_L = \qty{100}{\ohm}$. La captura A no cambia: $Z_{in}$ casi no depende de la carga.
\end{enumerate}

\paso{4}{Anote}
\begin{table}[H]
  \centering
  \renewcommand{\arraystretch}{1.15}
  \begin{tabularx}{\linewidth}{|>{\centering\arraybackslash}m{3.0cm}|*{3}{>{\centering\arraybackslash}X|}}
    \hline
    \textbf{Punto Q} & Teórico & Simulado & Medido \\ \hline
    $V_B$ & \hueco & & \\ \hline
    $V_E$ & \hueco & & \\ \hline
    $V_C$ & \hueco & & \\ \hline
    $I_E$ y $r_e$ & \hueco & & \\ \hline
  \end{tabularx}
  \caption{Experimento 1: punto de operación.}
  \label{tab:exp1dc}
\end{table}

\begin{table}[H]
  \centering
  \renewcommand{\arraystretch}{1.15}
  \begin{tabularx}{\linewidth}{|>{\centering\arraybackslash}m{3.0cm}|*{6}{>{\centering\arraybackslash}X|}}
    \hline
    & \multicolumn{3}{c|}{$\boldsymbol{R_L = \qty{10}{\kilo\ohm}}$}
    & \multicolumn{3}{c|}{$\boldsymbol{R_L = \qty{100}{\ohm}}$} \\
    \cline{2-7}
    \textbf{Parámetro} & Teórico & Simulado & Medido & Teórico & Simulado & Medido \\
    \hline
    $v_s$ (pico)      & \qty{50}{\milli\volt} & & & \qty{50}{\milli\volt} & & \\ \hline
    $v_{in}$ (pico)   & \hueco & & & & & \\ \hline
    $v_o$ (pico)      & \hueco & & & & & \\ \hline
    $Z_{in}$          & \hueco & & & & & \\ \hline
    $A_v$             & \hueco & & & & & \\ \hline
    $A_i$             & \hueco & & & & & \\ \hline
    $A_p$             & \hueco & & & & & \\ \hline
    $A_p$ (\unit{dB}) & \hueco & & & & & \\ \hline
  \end{tabularx}
  \caption{Experimento 1: ganancias con carga. $Z_{in}$ con~\eqref{eq:zin}; $A_i$ y $A_p$
           con~\eqref{eq:aiap}.}
  \label{tab:exp1ac}
\end{table}

\begin{table}[H]
  \centering
  \renewcommand{\arraystretch}{1.15}
  \begin{tabularx}{0.75\linewidth}{|>{\centering\arraybackslash}m{3.2cm}|*{3}{>{\centering\arraybackslash}X|}}
    \hline
    \textbf{Sin carga} & Teórico & Simulado & Medido \\ \hline
    $v_o$ (pico) & \hueco & & \\ \hline
    $A_v$        & \hueco & & \\ \hline
  \end{tabularx}
  \caption{Experimento 1: ganancia de voltaje sin carga.}
  \label{tab:exp1sin}
\end{table}

\noindent\begin{minipage}{\linewidth}
\noindent\textbf{Captura B con $R_L = \qty{10}{\kilo\ohm}$} (paso~3\ref{paso:carga}).
\cuadrorespuesta
\end{minipage}\par

%-----------------------------------------------------------------------
\subsection{Experimento 2: amplificador de voltaje en base común}
%-----------------------------------------------------------------------
Se arma sobre la misma placa del experimento~1, con el mismo punto Q. Solo cambian el terminal
por donde entra la señal y el que queda a tierra para la señal (tabla~\ref{tab:configuraciones}).

\paso{1}{Datos}
$V_{CC} = \qty{12}{\volt}$ (riel fijo); \texttt{W1}: senoidal de \qty{1}{\kilo\hertz} y
\qty{50}{\milli\volt} de amplitud; transistor BC548B; $R_s = \qty{100}{\ohm}$. La resistencia de
emisor queda entera sin desacoplar: $R_E = R_{E1} + R_{E2} = \qty{490}{\ohm}$. Se prueban dos
cargas: sin $R_L$ y $R_L = \qty{10}{\kilo\ohm}$. Los valores están en la
figura~\ref{fig:cb}.

\begin{figure}[H]
  \centering
  \begin{circuitikz}[scale=0.85, transform shape]
    \draw (5.6,0) node[npn](Q){};
    \node[left] at ($(Q.C)+(-0.1,0.1)$) {BC548B};
    \coordinate (gnd) at (0,-5.4);
    \coordinate (top) at (0,3.6);
    \coordinate (b) at ($(Q.B)+(-2.1,0)$);
    \draw (b |- top) to[R, l_=$R_1$, a^=\qty{10}{\kilo\ohm}] (b);
    \draw (b) to[R, l_=$R_2$, a^=\qty{1.8}{\kilo\ohm}] (b |- gnd);
    \draw (b) to[short, *-] (Q.B);
    \draw (b) -- ++(-2.2,0) coordinate(cb)
          to[C, l_={$C_B$\ \qty{100}{\micro\farad}}] (cb |- gnd);
    \coordinate (c) at ($(Q.C)+(0,0.6)$);
    \draw (c |- top) to[R, l_=$R_C$, a^=\qty{2.2}{\kilo\ohm}] (c) -- (Q.C);
    \draw (cb |- top) -- (c |- top) node[circ]{} -- ++(1.6,0) node[right]{$V_{CC} = \qty{+12}{\volt}$};
    \coordinate (e) at ($(Q.E)+(0,-1.0)$);
    \draw (Q.E) -- (e) to[R, l_=$R_{E1}$, a^=\qty{100}{\ohm}] ++(0,-1.8) coordinate(e2)
                       to[R, l_=$R_{E2}$, a^=\qty{390}{\ohm}] (e2 |- gnd);
    \draw (e) to[C, l={$C_{in}$\ \qty{100}{\micro\farad}}, *-] ++(2.6,0) coordinate(in)
              to[R, l_=$R_s$, a^=\qty{100}{\ohm}] ++(2.4,0) coordinate(s)
              to[sV, l_=\texttt{W1}] (s |- gnd);
    \node[above left] at (e) {$v_{in}$};
    \node[above] at (s) {$v_s$};
    \draw (c) to[C, l=$C_{out}$, a=\qty{10}{\micro\farad}, *-] ++(2.6,0) -- ++(3.2,0) coordinate(o)
              to[short, *-o] ++(1.0,0) node[right]{$v_o$};
    \draw (o) to[R, l=$R_L$] (o |- gnd);
    \draw (cb |- gnd) -- (o |- gnd);
    \draw (b |- gnd) node[ground]{};
  \end{circuitikz}
  \caption{Experimento 2: amplificador en base común. $C_B$ pone la base a tierra para la
           señal (franja hacia tierra) y la señal entra por el emisor. La franja de $C_{in}$
           va hacia $R_s$: el emisor está a unos \qty{1.1}{\volt}. $v_{in}$ se mide en el
           emisor.}
  \label{fig:cb}
\end{figure}

\paso{2}{Calcule (en casa)}
\begin{enumerate}[label=\alph*), leftmargin=*, itemsep=3pt]
  \item Explique por qué el punto Q es el mismo del experimento~1 (análisis DC: los
        capacitores son circuitos abiertos).
  \item $Z_{in}$ con~\eqref{eq:zincb} y el $v_{in}$ que llega al emisor:
        $v_{in} = v_s\,Z_{in}/(R_s + Z_{in})$.
  \item $A_v$ sin carga y con $R_L = \qty{10}{\kilo\ohm}$ con~\eqref{eq:avcb}, y la ganancia
        desde la fuente con~\eqref{eq:avscb}.
  \item $A_i$ y $A_p$ con $R_L = \qty{10}{\kilo\ohm}$ con~\eqref{eq:aiap}. Compruebe que $A_i$
        coincide con~\eqref{eq:aicb}.
  \item Simule el circuito con \qty{50}{\milli\volt} y con \qty{0.2}{\volt} de amplitud en
        \texttt{W1}.
\end{enumerate}

\paso{3}{Arme y mida (en el laboratorio)}
\begin{enumerate}[label=\alph*), ref=\alph*, leftmargin=*, itemsep=3pt]
  \item Sobre la placa del experimento~1: retire $C_E$ y $C_{in}$, cambie $R_s$ por
        \qty{100}{\ohm}, conecte $C_B$ de la base a tierra y lleve la señal al emisor a través
        del nuevo $C_{in}$ de \qty{100}{\micro\farad}.
  \item Con \texttt{W1} apagado, mida $V_B$, $V_E$ y $V_C$: deben ser los del experimento~1.
  \item \label{paso:cbmed}Encienda \texttt{W1} y haga las capturas A y B. En el emisor la señal
        es de pocos milivoltios sobre \qty{1.1}{\volt} de DC: ajuste el \emph{offset} del canal
        a \qty{-1.1}{\volt}, use una escala de \qtyrange{2}{5}{\milli\volt} por división y
        active \emph{Average}. Compruebe que $v_o$ está \textbf{en fase} con $v_{in}$.
  \item Repita la captura B con $R_L = \qty{10}{\kilo\ohm}$.
  \item \label{paso:cbdist}Sin $R_L$, suba \texttt{W1} a \qty{0.2}{\volt} y observe cómo se
        deforma $v_o$: un semiciclo sale más grande que el otro.
\end{enumerate}

\paso{4}{Anote}
\begin{table}[H]
  \centering
  \renewcommand{\arraystretch}{1.15}
  \begin{tabularx}{\linewidth}{|>{\centering\arraybackslash}m{3.0cm}|*{6}{>{\centering\arraybackslash}X|}}
    \hline
    & \multicolumn{3}{c|}{\textbf{Sin carga}}
    & \multicolumn{3}{c|}{$\boldsymbol{R_L = \qty{10}{\kilo\ohm}}$} \\
    \cline{2-7}
    \textbf{Parámetro} & Teórico & Simulado & Medido & Teórico & Simulado & Medido \\
    \hline
    $V_B$, $V_E$, $V_C$ & \hueco & & & \multicolumn{3}{c|}{---} \\ \hline
    $v_s$ (pico)      & \qty{50}{\milli\volt} & & & \qty{50}{\milli\volt} & & \\ \hline
    $v_{in}$ (pico)   & \hueco & & & & & \\ \hline
    $v_o$ (pico)      & \hueco & & & & & \\ \hline
    $Z_{in}$          & \hueco & & & & & \\ \hline
    $A_v = v_o/v_{in}$ & \hueco & & & & & \\ \hline
    $v_o/v_s$         & \hueco & & & & & \\ \hline
    $A_i$             & \multicolumn{3}{c|}{---} & & & \\ \hline
    $A_p$ y $A_p$ (\unit{dB}) & \multicolumn{3}{c|}{---} & & & \\ \hline
  \end{tabularx}
  \caption{Experimento 2: el amplificador en base común.}
  \label{tab:exp2}
\end{table}

\begin{table}[H]
  \centering
  \renewcommand{\arraystretch}{1.15}
  \begin{tabularx}{\linewidth}{|>{\centering\arraybackslash}m{4.6cm}|*{2}{>{\centering\arraybackslash}X|}}
    \hline
    \textbf{Medido, con $\boldsymbol{R_L = \qty{10}{\kilo\ohm}}$} & \textbf{Emisor común (exp.~1)}
      & \textbf{Base común (exp.~2)} \\ \hline
    $A_v$ (con su signo) & \hueco & \\ \hline
    $Z_{in}$             & \hueco & \\ \hline
    $A_i$                & \hueco & \\ \hline
    $A_p$                & \hueco & \\ \hline
  \end{tabularx}
  \caption{Experimento 2: las dos configuraciones con el mismo punto Q.}
  \label{tab:cecb}
\end{table}

\noindent\begin{minipage}{\linewidth}
\noindent\textbf{Captura B sin carga} (paso~3\ref{paso:cbmed}).
\cuadrorespuesta
\end{minipage}\par

\noindent\begin{minipage}{\linewidth}
\noindent\textbf{Captura de la deformación con \qty{0.2}{\volt}} (paso~3\ref{paso:cbdist}).
\cuadrorespuesta
\end{minipage}\par

%-----------------------------------------------------------------------
\subsection{Experimento 3: amplificador de corriente en colector común}
%-----------------------------------------------------------------------
Comienza la semana~2. El experimento~1 mostró que el emisor común no puede alimentar una carga
de \qty{100}{\ohm}. Un seguidor de emisor sí puede: no amplifica el voltaje, pero sí la
corriente. Para comprobarlo se comparan dos corrientes: la que entrega la fuente de señal,
$i_{in} = (v_s - v_{in})/R_s$, y la que recibe la carga, $i_o = v_o/R_L$.

\paso{1}{Datos}
$V_{CC} = \qty{12}{\volt}$ (riel fijo); \texttt{W1}: senoidal de \qty{1}{\kilo\hertz} y
\qty{0.4}{\volt} de amplitud; transistor 2N2222A; $R_L = \qty{100}{\ohm}$. Los valores están
en la figura~\ref{fig:cc}.

\begin{figure}[H]
  \centering
  \begin{circuitikz}[scale=0.85, transform shape]
    \draw (7.4,0) node[npn](Q){};
    \node[right] at (Q.east) {2N2222A};
    \coordinate (gnd) at (0,-4.3);
    \coordinate (top) at (0,2.6);
    \coordinate (b) at ($(Q.B)+(-1.0,0)$);
    % divisor de base
    \draw (b |- top) to[R, l_=$R_1$, a^=\qty{47}{\kilo\ohm}] (b);
    \draw (b) to[R, l_=$R_2$, a^=\qty{47}{\kilo\ohm}] (b |- gnd);
    \draw (b) to[short, *-] (Q.B);
    % entrada
    \coordinate (in) at ($(b)+(-2.4,0)$);
    \coordinate (s)  at ($(in)+(-2.6,0)$);
    \draw (in) to[C, l=$C_{in}$, a=\qty{10}{\micro\farad}] (b);
    \draw (s) to[R, l=$R_s$, a=\qty{10}{\kilo\ohm}, -*] (in);
    \node[above] at (in) {$v_{in}$};
    \draw (s |- gnd) to[sV, l=\texttt{W1}] (s);
    \node[above] at (s) {$v_s$};
    % colector directo a VCC
    \draw (Q.C) -- (Q.C |- top);
    \draw (b |- top) -- (Q.C |- top) node[circ]{} -- ++(3.6,0) node[right]{$V_{CC} = \qty{+12}{\volt}$};
    % emisor
    \coordinate (e) at ($(Q.E)+(0,-0.5)$);
    \draw (Q.E) -- (e) to[R, l_=$R_E$, a^=\qty{1}{\kilo\ohm}] (e |- gnd);
    \draw (e) to[C, l=$C_{out}$, a=\qty{100}{\micro\farad}, *-] ++(3.6,0) coordinate(o)
              to[R, l=$R_L$, a=\qty{100}{\ohm}, *-] (o |- gnd);
    \draw (o) to[short, -o] ++(1.0,0) node[right]{$v_o$};
    % tierra
    \draw (s |- gnd) -- (o |- gnd);
    \draw (Q.E |- gnd) node[ground]{};
  \end{circuitikz}
  \caption{Experimento 3: seguidor de emisor. La franja de $C_{out}$ va hacia $R_L$ (el
           emisor está a unos \qty{4.7}{\volt}).}
  \label{fig:cc}
\end{figure}

\paso{2}{Calcule (en casa)}
\begin{enumerate}[label=\alph*), leftmargin=*, itemsep=3pt]
  \item \textbf{Sin transistor:} si $R_L$ se conecta directamente detrás de $R_s$, ¿cuánto vale
        $v_o = v_s\,R_L/(R_s + R_L)$? Es lo que la fuente de señal puede dar sola.
  \item El punto Q ($V_{TH}$, $R_{TH}$, $I_E$, $V_E$) y $r_e$.
  \item $A_v$ con~\eqref{eq:avcc} y $Z_{in}$ con~\eqref{eq:zincc}; $i_{in}$, $i_o$, y $A_i$ y
        $A_p$ con~\eqref{eq:aiap}. Compárelos con los del emisor común con
        $R_L = \qty{100}{\ohm}$.
  \item El voltaje de salida a partir del cual se recorta el semiciclo negativo,
        con~\eqref{eq:recorte}.
  \item Calcule la ganancia de la \textbf{cascada} de la figura~\ref{fig:cascada}: la carga
        del emisor común es ahora la $Z_{in}$ del seguidor, así que
        $A_{v1} = -(R_C \pll Z_{in2})/(r_e + R_{E1})$, y la ganancia total es
        $A_v = A_{v1}\,A_{v2}$.
  \item Simule el seguidor con \qty{0.4}{\volt} y con \qty{1.2}{\volt} de amplitud en
        \texttt{W1}, y la cascada con \qty{20}{\milli\volt}.
\end{enumerate}

\paso{3}{Arme y mida (en el laboratorio)}
\begin{enumerate}[label=\alph*), ref=\alph*, leftmargin=*, itemsep=3pt]
  \item \textbf{Sin transistor:} conecte $R_L = \qty{100}{\ohm}$ directamente detrás de
        $R_s = \qty{10}{\kilo\ohm}$, con \texttt{W1} en \qty{0.4}{\volt}, y mida $v_o$.
  \item Arme el seguidor. Con \texttt{W1} apagado, mida $V_B$ y $V_E$.
  \item Encienda \texttt{W1} con \qty{0.4}{\volt} y haga las capturas A y B. Compruebe que
        $v_o$ está en fase con $v_{in}$ y casi del mismo tamaño, y compárelo con el $v_o$ sin
        transistor.
  \item \label{paso:recorte}Suba la amplitud de \texttt{W1} hasta que el semiciclo negativo de
        $v_o$ se aplane. Capture la forma de onda y mida el voltaje donde se recorta.
  \item \label{paso:cascada}Arme la cascada de la figura~\ref{fig:cascada}. Vuelva a armar el
        emisor común del experimento~1 (con $C_E$ y sin $C_B$), quítele $R_s$ y $R_L$ y una su
        $C_{out}$ con la base del seguidor
        (en lugar de $C_{in}$ y $R_s$ del seguidor). Ponga \texttt{W1} en
        \qty{20}{\milli\volt} y mida $v_o$ en los \qty{100}{\ohm}.
\end{enumerate}

\begin{figure}[H]
  \centering
  \begin{tikzpicture}[>=Stealth, every node/.style={font=\small},
      bloque/.style={draw, thick, fill=black!5, minimum height=1.2cm, align=center,
                     text width=2.8cm}]
    \node (w1) [draw, circle, minimum size=0.9cm] {\texttt{W1}};
    \node (ce) [bloque, right=1.4cm of w1] {Emisor común\\(exp.~1, sin $R_L$)};
    \node (cc) [bloque, right=1.4cm of ce] {Seguidor de emisor\\(exp.~3, sin $R_s$)};
    \node (rl) [draw, thick, minimum width=0.6cm, minimum height=1.2cm,
                right=1.4cm of cc] {\rotatebox{90}{\qty{100}{\ohm}}};
    \draw[->] (w1) -- node[above]{$v_{in}$} (ce);
    \draw[->] (ce) -- node[above]{$C_{out}$} (cc);
    \draw[->] (cc) -- node[above]{$v_o$} (rl);
  \end{tikzpicture}
  \caption{Experimento 3: cascada. El emisor común da la ganancia de voltaje y el seguidor la
           corriente que necesita la carga de \qty{100}{\ohm}.}
  \label{fig:cascada}
\end{figure}

\paso{4}{Anote}
\begin{table}[H]
  \centering
  \renewcommand{\arraystretch}{1.15}
  \begin{tabularx}{\linewidth}{|>{\centering\arraybackslash}m{4.4cm}|*{3}{>{\centering\arraybackslash}X|}}
    \hline
    \textbf{Parámetro} & Teórico & Simulado & Medido \\ \hline
    $v_o$ sin transistor   & \hueco & & \\ \hline
    $V_B$ y $V_E$          & \hueco & & \\ \hline
    $I_E$ y $r_e$          & \hueco & & \\ \hline
    $v_s$ (pico)           & \qty{0.4}{\volt} & & \\ \hline
    $v_{in}$ (pico)        & \hueco & & \\ \hline
    $v_o$ (pico)           & \hueco & & \\ \hline
    $Z_{in}$               & \hueco & & \\ \hline
    $A_v$                  & \hueco & & \\ \hline
    $i_{in} = (v_s - v_{in})/R_s$ & \hueco & & \\ \hline
    $i_o = v_o/R_L$        & \hueco & & \\ \hline
    $A_i = i_o/i_{in}$     & \hueco & & \\ \hline
    $A_p$ y $A_p$ (\unit{dB}) & \hueco & & \\ \hline
    $v_o$ donde se recorta & \hueco & & \\ \hline
  \end{tabularx}
  \caption{Experimento 3: el seguidor de emisor con $R_L = \qty{100}{\ohm}$.}
  \label{tab:exp3}
\end{table}

\begin{table}[H]
  \centering
  \renewcommand{\arraystretch}{1.15}
  \begin{tabularx}{\linewidth}{|>{\centering\arraybackslash}m{5.6cm}|*{3}{>{\centering\arraybackslash}X|}}
    \hline
    \textbf{$A_v$ en \qty{100}{\ohm}} & Teórico & Simulado & Medido \\ \hline
    Emisor común solo (tabla~\ref{tab:exp1ac}) & \hueco & & \\ \hline
    Emisor común + seguidor (cascada)          & \hueco & & \\ \hline
  \end{tabularx}
  \caption{Experimento 3: la misma carga con y sin seguidor.}
  \label{tab:cascada}
\end{table}

\noindent\begin{minipage}{\linewidth}
\noindent\textbf{Captura del recorte del seguidor} (paso~3\ref{paso:recorte}).
\cuadrorespuesta
\end{minipage}\par

%-----------------------------------------------------------------------
\subsection{Experimento 4: amplificador de potencia push-pull}
%-----------------------------------------------------------------------
Dos seguidores complementarios entregan potencia a una carga de \qty{47}{\ohm}. Primero se
arma la clase B para ver la distorsión de cruce y luego se corrige con la clase AB.

\paso{1}{Datos}
Fuentes programables de la tarjeta: $V_+ = \qty{+5}{\volt}$ y $V_- = \qty{-5}{\volt}$;
\texttt{W1}: senoidal de \qty{1}{\kilo\hertz} y \qty{3}{\volt} de amplitud; transistores
2N2222A (NPN) y 2N2907A (PNP); $R_L = \qty{47}{\ohm}$, \qty{0.5}{\watt};
$R_s = \qty{100}{\ohm}$.

\begin{figure}[H]
  \centering
  \begin{subfigure}[b]{0.46\linewidth}
    \centering
    \begin{circuitikz}[scale=0.78, transform shape]
      \draw (5,1.8) node[npn](Q1){};
      \draw (5,-1.8) node[pnp](Q2){};
      \node[right] at (Q1.east) {2N2222A};
      \node[right] at (Q2.east) {2N2907A};
      \draw (Q1.E) -- (Q2.E);
      \coordinate (o) at ($(Q1.E)!0.5!(Q2.E)$);
      \coordinate (bb) at ($(Q1.B)!0.5!(Q2.B)+(-0.8,0)$);
      \draw (Q1.B) -| (bb) |- (Q2.B);
      \draw (Q1.C) -- ++(0,0.5) node[above]{$V_+ = \qty{+5}{\volt}$};
      \draw (Q2.C) -- ++(0,-0.5) node[below]{$V_- = \qty{-5}{\volt}$};
      \draw (o) to[short, *-] ++(2.1,0) coordinate(ro)
                to[R, l=$R_L$] ++(0,-2.2) node[ground]{};
      \draw (ro) to[short, *-o] ++(0.8,0) node[above]{$v_o$};
      \coordinate (s) at ($(bb)+(-2.2,0)$);
      \draw (s) to[R, l=$R_s$, -*] (bb);
      \node[below left] at (bb) {$v_{in}$};
      \draw (s) to[sV, l_=\texttt{W1}] ++(0,-2.2) node[ground]{};
    \end{circuitikz}
    \caption{Clase B.}
    \label{fig:claseb}
  \end{subfigure}
  \hfill
  \begin{subfigure}[b]{0.50\linewidth}
    \centering
    \begin{circuitikz}[scale=0.78, transform shape]
      \draw (5,1.8) node[npn](Q1){};
      \draw (5,-1.8) node[pnp](Q2){};
      \node[right] at (Q1.east) {2N2222A};
      \node[right] at (Q2.east) {2N2907A};
      \draw (Q1.E) -- (Q2.E);
      \coordinate (o) at ($(Q1.E)!0.5!(Q2.E)$);
      \coordinate (b1) at ($(Q1.B)+(-0.9,0)$);
      \coordinate (b2) at ($(Q2.B)+(-0.9,0)$);
      \coordinate (m)  at ($(b1)!0.5!(b2)$);
      \draw (Q1.B) -- (b1);
      \draw (Q2.B) -- (b2);
      \draw (b1) to[D, l=$D_1$, *-*] (m) to[D, l=$D_2$, -*] (b2);
      \draw (Q1.C) |- ($(b1)+(0,2.0)$) coordinate(tp);
      \draw (tp) to[R, l_=$R_{B1}$, a^=\qty{1}{\kilo\ohm}] (b1);
      \draw (Q2.C) |- ($(b2)+(0,-2.0)$) coordinate(bt);
      \draw (b2) to[R, l_=$R_{B2}$, a^=\qty{1}{\kilo\ohm}] (bt);
      \node[above] at ($(Q1.C |- tp)$) {$V_+$};
      \node[below] at ($(Q2.C |- bt)$) {$V_-$};
      \draw (o) to[short, *-] ++(2.1,0) coordinate(ro)
                to[R, l=$R_L$] ++(0,-2.2) node[ground]{};
      \draw (ro) to[short, *-o] ++(0.8,0) node[above]{$v_o$};
      \coordinate (in) at ($(m)+(-2.4,0)$);
      \draw (in) to[C, l=$C_{in}$, a=\qty{47}{\micro\farad}] (m);
      \coordinate (s) at ($(in)+(-1.8,0)$);
      \draw (s) to[R, l=$R_s$, -*] (in);
      \node[below] at (in) {$v_{in}$};
      \draw (s) to[sV, l_=\texttt{W1}] ++(0,-2.2) node[ground]{};
    \end{circuitikz}
    \caption{Clase AB.}
    \label{fig:claseab}
  \end{subfigure}
  \caption{Experimento 4: etapa de potencia push-pull. En la clase AB los diodos 1N4148 (con
           la franja hacia $V_-$) mantienen a cada transistor al borde de conducir.}
  \label{fig:pushpull}
\end{figure}

\paso{2}{Calcule (en casa)}
\begin{enumerate}[label=\alph*), leftmargin=*, itemsep=3pt]
  \item Para la clase B, con $V_p \approx \qty{3}{\volt} - \qty{0.7}{\volt}$: $P_{out}$
        con~\eqref{eq:pout}, $P_{DC}$ con~\eqref{eq:pdc}, $\eta$ con~\eqref{eq:eta} y la
        potencia máxima que disipa cada transistor. Compruebe que no pasa de la que admite el
        encapsulado TO-92.
  \item Para la clase AB: la corriente por los diodos, $(V_+ - V_- - 2V_D)/(R_{B1} + R_{B2})$,
        y la potencia que agrega a $P_{DC}$. Estime $Z_{in}$ (casi toda es
        $R_{B1} \pll R_{B2}$).
  \item Simule las dos clases. Anote $v_o$ pico y eficaz (rms), las corrientes promedio de
        $V_+$ y $V_-$ y, en la clase AB, $v_{in}$ y la corriente eficaz por $R_s$.
\end{enumerate}

\paso{3}{Arme y mida (en el laboratorio)}
\begin{enumerate}[label=\alph*), ref=\alph*, leftmargin=*, itemsep=3pt]
  \item Programe $V_+ = \qty{5}{\volt}$ y $V_- = \qty{-5}{\volt}$ en WaveForms
        (\emph{Supplies}) y arme la clase B. Ponga el multímetro, en la escala de
        \unit{\milli\ampere} DC, en serie con $V_+$; necesitará medir también $I_-$.
  \item \label{paso:cruce}Con \texttt{W1} en \qty{3}{\volt}, capture $v_{in}$ (CH1) y $v_o$
        (CH2). Haga zoom sobre un cruce por cero y mida el ancho de la muesca. Anote
        $v_o$ pico, $v_o$ rms e $I_+$; luego mueva el multímetro a $V_-$ y anote $I_-$.
  \item Arme la clase AB y repita el paso anterior. Añada la captura A (CH1 en $v_s$, CH2 en
        $v_{in}$, en valores rms) para calcular la corriente de entrada,
        $i_{in(\text{rms})} = (v_{s} - v_{in})/R_s$.
  \item Toque los transistores: deben estar apenas tibios. Si se calientan, apague y revise
        el montaje.
\end{enumerate}

\paso{4}{Anote}
\begin{table}[H]
  \centering
  \renewcommand{\arraystretch}{1.15}
  \begin{tabularx}{\linewidth}{|>{\centering\arraybackslash}m{3.6cm}|*{6}{>{\centering\arraybackslash}X|}}
    \hline
    & \multicolumn{3}{c|}{\textbf{Clase B}} & \multicolumn{3}{c|}{\textbf{Clase AB}} \\
    \cline{2-7}
    \textbf{Parámetro} & Teórico & Simulado & Medido & Teórico & Simulado & Medido \\
    \hline
    $v_o$ pico                       & \hueco & & & & & \\ \hline
    $v_o$ rms                        & \hueco & & & & & \\ \hline
    $A_v = v_{o}/v_{in}$ (rms)       & \hueco & & & & & \\ \hline
    $I_+$ e $I_-$ (promedio)         & \hueco & & & & & \\ \hline
    $P_{out} = v_{o(\text{rms})}^2/R_L$ & \hueco & & & & & \\ \hline
    $P_{DC}$                         & \hueco & & & & & \\ \hline
    $\eta = P_{out}/P_{DC}$          & \hueco & & & & & \\ \hline
    $i_{in}$ rms                     & \multicolumn{3}{c|}{---} & & & \\ \hline
    $P_{in} = v_{in}\,i_{in}$ (rms)  & \multicolumn{3}{c|}{---} & & & \\ \hline
    $A_p = P_{out}/P_{in}$           & \multicolumn{3}{c|}{---} & & & \\ \hline
  \end{tabularx}
  \caption{Experimento 4: la etapa de potencia. En la clase B la corriente de entrada es tan
           pequeña que no se puede medir con $R_s = \qty{100}{\ohm}$.}
  \label{tab:exp4}
\end{table}

\noindent\begin{minipage}{\linewidth}
\noindent\textbf{Captura de la distorsión de cruce en clase B} (paso~3\ref{paso:cruce}).
\cuadrorespuesta
\end{minipage}\par

\noindent\begin{minipage}{\linewidth}
\noindent\textbf{Captura de la salida en clase AB.}
\cuadrorespuesta
\end{minipage}\par

\FloatBarrier
%=======================================================================
\section{Informe final}
%=======================================================================
\begin{enumerate}[label=\textbf{\arabic*.}, leftmargin=*, labelsep=0.6em]
  \item Adjunte los cálculos y las simulaciones de los cuatro experimentos.
  \item Compare los valores teóricos, simulados y medidos y explique las diferencias más
        grandes.
  \item Con las tablas~\ref{tab:cecb}, \ref{tab:exp3} y~\ref{tab:cascada}, explique qué
        amplifica cada configuración (emisor común, base común y colector común) y por qué la
        cascada funciona mejor que cada etapa por separado.
  \item Complete la tabla~\ref{tab:comparacion} con los valores \emph{medidos} de $A_v$, $A_i$
        y $Z_{in}$ de cada experimento.
\end{enumerate}

\FloatBarrier
%=======================================================================
\section{Cuestionario}
%=======================================================================
\begin{enumerate}[label=\textbf{\arabic*.}, leftmargin=*, labelsep=0.6em]
\item En el experimento 1, ¿cuánto valdrían $A_v$ y $Z_{in}$ (sin carga) si se retira el
capacitor $C_E$? ¿Qué se gana y qué se pierde al dejar sin desacoplar toda la resistencia de
emisor?
\lineasresp{3}

\item En el experimento 2, ¿cuánto valdría la ganancia desde la fuente, $v_o/v_s$ (sin carga),
si $R_s$ fuera de \qty{1}{\kilo\ohm}? ¿Por qué la base común necesita una fuente de señal de
impedancia muy baja, y por qué entonces se usa en altas frecuencias?
\lineasresp{3}

\item El seguidor de emisor tiene $A_v < 1$. ¿De dónde sale entonces la potencia extra que
recibe la carga? ¿Hasta qué amplitud de salida funciona sin recortar con
$R_L = \qty{100}{\ohm}$, y cómo cambiaría ese límite si $R_E$ fuera de \qty{470}{\ohm}?
\lineasresp{3}

\item ¿Cuál sería el rendimiento de la clase B si la salida llegara a $V_p = \qty{4.3}{\volt}$?
¿Por qué el rendimiento medido en la clase AB es menor que el de la clase B, si la señal sale
mejor?
\lineasresp{3}
\end{enumerate}

\FloatBarrier
%=======================================================================
\section{Conclusiones}
%=======================================================================
\begin{itemize}
    \item Redacte como mínimo 05 conclusiones \textbf{cualitativas} y 05
          \textbf{cuantitativas}.
\end{itemize}

\noindent\textbf{Cualitativas}
\lineasresp{10}

\needspace{6\baselineskip}
\noindent\textbf{Cuantitativas}
\lineasresp{10}

\end{document}
